\ea\label{ex:text19-1}
GOROTO\\
\gll goroto\\
     turtle\\
\glt `TURTLE'
\z

\ea\label{ex:text19-2}
A goroto a iana a beera sana.\\
\gll a	goroto	a	iana	a	beera	sana\\
     \textsc{art}2.\textsc{sg}	turtle	\textsc{art}2.\textsc{sg}	swimming.animal	\textsc{art}2.\textsc{sg}	big	very\\
\glt `The turtle is a very big swimming animal.'
\z

\ea\label{ex:text19-3}
Bean tara e, ean repaa tara voen:\\
\gll be=an	tara	e	e-an	re=paa	tara	vo=en\\
     when=2\textsc{sg}.\textsc{pron}	see	3\textsc{sg}.\textsc{pron}	\textsc{h}-2\textsc{sg}.\textsc{pron}	\textsc{consec}=\textsc{tm}3	see	like=\textsc{dem}3\\
\glt `When you see it, you see:'
\z

\ea\label{ex:text19-4}
A goroto na tei me nana bono kahoo, o dorana o pakapaka.\\
\gll a	goroto	na	tei	me	nana	bono	kahoo	o	dorana	o	pakapaka\\
     \textsc{art}2.\textsc{sg}	turtle	\textsc{tm}2	be.there	with	3\textsc{sg}.\textsc{ipfv}	\textsc{obj}.\textsc{art}3.\textsc{sg}	head	\textsc{art}3.\textsc{sg}	four	\textsc{art}3.\textsc{sg}	flipper\\
\glt `The turtle has a head, (and) four flippers.'
\z

\ea\label{ex:text19-5}
O kapa vaa toonunae o kikis, o kapa vaa domonae saka kikis haana.\\
\gll o	kapa	vaa	toonu-nae	o	kikis	o	kapa	vaa	domo-nae	saka	kikis	haa=na\\
     \textsc{art}3.\textsc{sg}	shell	\textsc{lnk}	back-3\textsc{sg}.\textsc{poss}.\textsc{pron}	\textsc{art}3.\textsc{sg}	hard	\textsc{art}3.\textsc{sg}	shell	\textsc{lnk}	chest-3\textsc{sg}.\textsc{poss}.\textsc{pron}	\textsc{neg}1	hard	\textsc{neg}2=3\textsc{sg}.\textsc{ipfv}\\
\glt `The shell on its back is hard, the shell on its chest is not hard.'
\z

\ea\label{ex:text19-6}
Teo dorana o pakapaka, o buaku vaa kahoo nae na dao raara bono huahua.\\
\gll te=o	dorana	o	pakapaka	o	buaku	vaa	kahoo-	nae	na	dao	ra=ara	bono	hua{$\sim$}hua\\
     \textsc{prep}=\textsc{art}3.\textsc{sg}	four	\textsc{art}3.\textsc{sg}	flipper	\textsc{art}3.\textsc{sg}	two	\textsc{lnk}	head-	3\textsc{sg}.\textsc{poss}.\textsc{pron}	\textsc{tm}2	call	1\textsc{pl}.\textsc{in}.\textsc{ipfv}=1\textsc{pl}.\textsc{in}.\textsc{pron}	\textsc{obj}.\textsc{art}3.\textsc{sg}	\textsc{red}{$\sim$}paddle\\
\glt `As for the four flippers, the two flippers near its head, we call paddles.'
\z

\newpage
\ea\label{ex:text19-7}
A goroto na huahua vakikis ni nana bono buo pakapaka bona.\\
\gll a	goroto	na	hua{$\sim$}hua	va-kikis	ni	nana	bono	buo	pakapaka	bona\\
     \textsc{art}2.\textsc{sg}	turtle	\textsc{tm}2	\textsc{red}{$\sim$}paddle	\textsc{advr}-strong	\textsc{appl}	3\textsc{sg}.\textsc{ipfv}	\textsc{obj}.\textsc{art}3.\textsc{sg}	two3.\textsc{sg}	flipper	\textsc{ana}.\textsc{sg}\\
\glt `The turtle strongly paddles with these two flippers.'
\z

\ea\label{ex:text19-8}
O buo pakapaka vaa murinae to vakarovo bata ni nae bara na hua ni nae mee.\\
\gll o	buo	pakapaka	vaa	muri-nae	to	vakarovo	bata	ni	na=e	bara	na	hua	ni	na=e	mee\\
     \textsc{art}3.\textsc{sg}	two3.\textsc{sg}	flipper	\textsc{lnk}	back-3\textsc{sg}.\textsc{poss}.\textsc{pron}	\textsc{rel}	steer	\textsc{cont}	\textsc{appl}	3\textsc{sg}.\textsc{ipfv}=3\textsc{sg}.\textsc{pron}	and2	\textsc{tm}2	paddle	\textsc{appl}	3\textsc{sg}.\textsc{ipfv}=3\textsc{sg}.\textsc{pron}	also\\
\glt `The two hind flippers it uses for steering and also for paddling.'\footnote{lit. `The two flippers at its back (are the ones) that it steers with.'}
\z

\ea\label{ex:text19-9}
A goroto na vaavuriha nana mee be tavusu kahi bona tahii, ore paa ruuvu gunaha vahaa.\\
\gll a	goroto	na	vaavuriha	nana	mee	be	tavusu	kahi	bona	tahii	o=re	paa	ruuvu	gunaha	vahaa\\
     \textsc{art}2.\textsc{sg}	turtle	\textsc{tm}2	breathe	3\textsc{sg}.\textsc{ipfv}	also	when	come.out	from	\textsc{obj}.\textsc{art}2.\textsc{sg}	salt.water	3.\textsc{pron}=\textsc{consec}	\textsc{tm}3	dive	go.down	again\\
\glt `The turtle also breathes when it comes out of the water, then it dives down again.'
\z

\ea\label{ex:text19-10}
A goroto vai a moon be tau vahuhu, ore paa gunaha vo teo kakaamuru.\\
\gll a	goroto	vai	a	moon	be	tau	vahuhu	o=re	paa	gunaha	vo	te=o	kakaamuru\\
     \textsc{art}2.\textsc{sg}	turtle	\textsc{dem}6	\textsc{art}2.\textsc{sg}	female	when	\textsc{tm}5	lay.eggs	3.\textsc{pron}=\textsc{consec}	\textsc{tm}3	go.ashore	\textsc{goal}	\textsc{prep}=\textsc{art}3.\textsc{sg}	white.sand\\
\glt `When the female turtle is about to lay eggs, it goes ashore to the white sand.'
\z

\ea\label{ex:text19-11}
Na kasi nana bona govee kasuana\\
\gll na	kasi	nana	bona	govee	kasuana\\
     \textsc{tm}2	dig	3\textsc{sg}.\textsc{ipfv}	\textsc{obj}.\textsc{art}2.\textsc{sg}	hole	ground\\
\glt `It digs a hole in the ground,'
\z

\newpage
\ea\label{ex:text19-12}
ore paa vahuhu komana govee\\
\gll o=re	paa	vahuhu	koma-n=a	govee\\
     3.\textsc{pron}=\textsc{consec}	\textsc{tm}3	lay.eggs	inside-3\textsc{sg}.\textsc{poss}=\textsc{art}2.\textsc{sg}	hole\\
\glt `then it lays its eggs into the hole'
\z

\ea\label{ex:text19-13}
ore paa pae vahaa bona\\
\gll o=re	paa	pae	vahaa	bona\\
     3.\textsc{pron}=\textsc{consec}	\textsc{tm}3	bury	again	4\textsc{sg}.\textsc{pron}\\
\glt `and fills the hole up again.'
\z

\ea\label{ex:text19-14}
Evehee bea goroto pae vahaa bona govee, na varovarova nae,\\
\gll evehee	be=a	goroto	pae	vahaa	bona	govee	na	va-rova{$\sim$}rova	na=e\\
     but	when=\textsc{art}2.\textsc{sg}	turtle	fill.hole	again	\textsc{obj}.\textsc{art}2.\textsc{sg}	hole	\textsc{tm}2	\textsc{mult}-\textsc{red}{$\sim$}disappear	3\textsc{sg}.\textsc{ipfv}=3\textsc{sg}.\textsc{pron}\\
\glt `But when the turtle fills up the hole again, it becomes invisable,'
\z

\ea\label{ex:text19-15}
a aba ore saka antee haa bona tea vurahe vakahu.\\
\gll a	aba	o=re	saka	antee	haa	bona	tea	vurahe	vakahu\\
     \textsc{art}2.\textsc{sg}	person	3.\textsc{pron}=\textsc{consec}	\textsc{neg}1	can	\textsc{neg}2	4\textsc{sg}.\textsc{pron}	\textsc{compl}1	find	quickly\\
\glt `so that one cannot quickly find it.'
\z

\ea\label{ex:text19-16}
A rupi a bero kurusu to vahuhu nana a goroto.\\
\gll a	rupi	a	bero	kurusu	to	vahuhu	nana	a	goroto\\
     \textsc{art}3.\textsc{pl}	egg	\textsc{art}3.\textsc{pl}	many	very	\textsc{rel}	lay.eggs	3\textsc{sg}.\textsc{ipfv}	\textsc{art}2.\textsc{sg}	turtle\\
\glt `The turtle lays very many eggs.'\footnote{lit. `(It is) many eggs indeed that the turtle lays.'}
\z

\ea\label{ex:text19-17}
Toro oha nana bono taonim saavun.\\
\gll toro	oha	nana	bono	taonim	saavun\\
     \textsc{tm}7	exceed	3\textsc{sg}.\textsc{ipfv}	\textsc{obj}.\textsc{art}3.\textsc{sg}	five	ten\\
\glt `Must be more than fifty.'
\z

\newpage
\ea\label{ex:text19-18}
Bea peha aba naonao tea ruuvu goroto, na nao vatotobin vo nana irihi.\\
\gll be=a	peha	aba	nao{$\sim$}nao	tea	ruuvu	goroto	na	nao	va-totobin	vo	nana	irihi\\
     when=\textsc{art}2.\textsc{sg}	\textsc{indef}2.\textsc{sg}	person	\textsc{red}{$\sim$}go	\textsc{compl}1	dive	turtle	\textsc{tm}2	go	\textsc{advr}-straight	\textsc{goal}	3\textsc{sg}.\textsc{ipfv}	edge.of.reef\\
\glt `When someone goes diving for turtles, he goes straight to the edge of the reef.'
\z

\ea\label{ex:text19-19}
O goroto na teitei rori irihi tea maa avaava vasu ei teo hum namana.\\
\gll o	goroto	na	tei{$\sim$}tei	rori	irihi	te=a	maa	avaava	vasu	ei	te=o	hum	namana\\
     \textsc{art}2.\textsc{pl}	turtle	\textsc{tm}2	\textsc{red}{$\sim$}be.there	3\textsc{pl}.\textsc{ipfv}	edge.of.reef	\textsc{prep}=\textsc{art}2.\textsc{sg}	\textsc{plm}	cave	stone	\textsc{dem}2	\textsc{prep}=\textsc{art}3.\textsc{sg}	place	ocean\\
\glt `The turtles stay at the edge of the reef in the caves of the rocks at the side of the ocean.'
\z

\ea\label{ex:text19-20}
Bea tahii masi, ore paa vurahe vakahu koa bono goroto,\\
\gll be=a	tahii	masi	o=re	paa	vurahe	vakahu	koa	bono	goroto\\
     when=\textsc{art}2.\textsc{sg}	sea	be.at.low.tide	3.\textsc{pron}=\textsc{consec}	\textsc{tm}3	find	quickly	just	\textsc{obj}.\textsc{art}2.\textsc{pl}	turtle\\
\glt `When the sea is at low tide (so that the reef is visable), then they can quickly find the turtles,'
\z

\ea\label{ex:text19-21}
beori ani rori bono vasukare bara tea maa avaava vasu ae tea maa vaagoroto.\\
\gll be=ori	ani	rori	bono	vasukare	bara	te=a	maa	avaava	vasu	ae	te=a	maa	vaagoroto\\
     when=3\textsc{pl}.\textsc{pron}	eat	3\textsc{pl}.\textsc{ipfv}	\textsc{obj}.\textsc{art}3.\textsc{sg}	vasukare.coral	and2	\textsc{prep}=\textsc{art}2.\textsc{sg}	\textsc{plm}	cave	stone	and1	\textsc{prep}=\textsc{art}2.\textsc{sg}	\textsc{plm}	turtle.tunnel\\
\glt `when they are eating \textit{vasukare} corals and stay in the rock caves and in the turtle tunnels.'
\z

\ea\label{ex:text19-22}
HAVEE TO GONO KAAKOO MI RIORI BONA GOROTO\\
\gll havee	to	gono	kaakoo	mi	ri=ori	bona	goroto\\
     how	\textsc{rel}	get	move	for	3\textsc{pl}.\textsc{ipfv}=3\textsc{pl}.\textsc{pron}	\textsc{obj}.\textsc{art}2.\textsc{sg}	turtle\\
\glt `HOW THEY CATCH A TURTLE'
\z

\newpage
\ea\label{ex:text19-23}
A goroto na ato riori tea kave,\\
\gll a	goroto	na	ato	ri=ori	te=a	kave\\
     \textsc{art}2.\textsc{sg}	turtle	\textsc{tm}2	get.hold.of	3\textsc{pl}.\textsc{ipfv}=3\textsc{pl}.\textsc{pron}	\textsc{prep}=\textsc{art}2.\textsc{sg}	net\\
\glt `They get hold of the turtle with the net,'\footnote{lit. `The turtle they get hold of with the net.'}
\z


\ea\label{ex:text19-24}
ei beori nao tea veo goroto.\\
\gll ei	be=ori	nao	tea	veo	goroto\\
     \textsc{dem}2	when=3\textsc{pl}.\textsc{pron}	go	\textsc{compl}1	net	turtle\\
\glt `that is when they go turtle netting.'\footnote{lit. `this (means) ...'}
\z

\ea\label{ex:text19-25}
A kave bona a beera sana.\\
\gll a	kave	bona	a	beera	sana\\
     \textsc{art}2.\textsc{sg}	net	\textsc{ana}.\textsc{sg}	\textsc{art}2.\textsc{sg}	big	very\\
\glt `This net is very big.'
\z

\ea\label{ex:text19-26}
Na antee nana tea kana bono buaku saavun goroto.\\
\gll na	antee	nana	tea	kana	bono	buaku	saavun	goroto\\
     \textsc{tm}2	can	3\textsc{sg}.\textsc{ipfv}	\textsc{compl}1	catch.in.net	\textsc{obj}.\textsc{art}3.\textsc{sg}	two	ten	turtle\\
\glt `It can catch twenty turtles.'
\z

\ea\label{ex:text19-27}
Amaa peha maa vuri to naonao riori tea ruuvu goroto.\\
\gll a=maa	peha	maa	vuri	to	nao{$\sim$}nao	ri=ori	tea	ruuvu	goroto\\
     \textsc{art}2.\textsc{sg}=\textsc{plm}	\textsc{indef}2.\textsc{sg}	\textsc{plm}	time	\textsc{rel}	\textsc{red}{$\sim$}go	3\textsc{pl}.\textsc{ipfv}=3\textsc{pl}.\textsc{pron}	\textsc{compl}1	dive	turtle\\
\glt `Sometimes they go diving for turles.'\footnote{lit. `(There are) some times that they go dive (for) turtles.'}
\z

\ea\label{ex:text19-28}
A goroto mee na kapa riori bea goroto gunaha nao tea vahuhu kasuana.\\
\gll a	goroto	mee	na	kapa	ri=ori	be=a	goroto	gunaha	nao	tea	vahuhu	kasuana\\
     \textsc{art}2.\textsc{sg}	turtle	also	\textsc{tm}2	lift	3\textsc{pl}.\textsc{ipfv}=3\textsc{pl}.\textsc{pron}	when=\textsc{art}2.\textsc{sg}	turtle	go.ashore	thither	\textsc{compl}1	lay.eggs	sand\\
\glt `They also catch a turtle when the turtle goes ashore to lay its eggs in the sand.'
\z

\ea\label{ex:text19-29}
Tea maa vuri sunuu, a goroto mee na vurahe riori ei be masi pina ni tea maa kopua.\\
\gll te=a	maa	vuri	sunuu	a	goroto	mee	na	vurahe	ri=ori	ei	be	masi	pina	ni	te=a	maa	koopua\\
     \textsc{prep}=\textsc{art}2.\textsc{sg}	\textsc{plm}	time	fish.with.torches	\textsc{art}2.\textsc{sg}	turtle	also	\textsc{tm}2	find	3\textsc{pl}.\textsc{ipfv}=3\textsc{pl}.\textsc{pron}	\textsc{dem}2	when	be.low.tide	blocking	\textsc{appl}	\textsc{prep}=\textsc{art}2.\textsc{sg}	\textsc{plm}	lagoon\\
\glt `When they go fishing with torches, they also find the turtle when it gets trapped in the lagoons at low tide.'\footnote{lit. `At the times (of) fishing with torches, they also find the turtle, this (is) when low tide blocking (it) in the lagoons.'}
\z

\ea\label{ex:text19-30}
HAVEE TO KAKU GOROTO VAAVAHA VO RIA ABANA VAA TEAPU\\
\gll havee	to	kaku	goroto	vaavaha	vo	ri=a	abana	vaa	teapu\\
     what	\textsc{rel}	butcher	turtle	different	like	3\textsc{pl}.\textsc{ipfv}=\textsc{art}2.\textsc{sg}	people	\textsc{lnk}	Teop.region\\
\glt `HOW TURTLES ARE BUTCHERED BY THE TEOP PEOPLE'\footnote{lit. `What (is it) like that the men of the Teop region butcher turtles differently.'}
\z

\ea\label{ex:text19-31}
Nabunuu roho, a kaku goroto, amaa abana bebeera koa to pakupaku roho bari.\\
\gll nabunuu	roho	a	kaku	goroto	a=maa	abana	be{$\sim$}beera	koa	to	paku{$\sim$}paku	roho	bari\\
     in.the.past	before	\textsc{art}2.\textsc{sg}	butcher	turtle	\textsc{art}2.\textsc{sg}=\textsc{plm}	men	\textsc{red}{$\sim$}big	only	\textsc{rel}	\textsc{red}{$\sim$}do	before	4\textsc{pl}.\textsc{pron}\\
\glt `In former times the butchering of turtles was only done by the big men.'\footnote{lit. `In former times, the butchering (of) turtles, only big men (were the ones) who did it.'}
\z

\ea\label{ex:text19-32}
Ean saka antee haa nom tea tara ta peha beiko ge ta peha visoasi teo hum vai to kaku goroto riori.\\
\gll e-an	saka	antee	haa	nom	tea	tara	ta	peha	beiko	ge	ta	peha	visoasi	te=o	hum	vai	to	kaku	goroto	ri=ori\\
     \textsc{h}-2\textsc{sg}.\textsc{pron}	\textsc{neg}1	can	\textsc{neg}2	\textsc{ipfv}	\textsc{compl}1	see	\textsc{nspec}2.\textsc{sg}	one2.\textsc{sg}	child	or	\textsc{nspec}2.\textsc{sg}	one2.\textsc{sg}	young.man	\textsc{prep}=\textsc{art}3.\textsc{sg}	place	\textsc{dem}6	\textsc{rel}	butcher	turtle	3\textsc{pl}.\textsc{ipfv}=3\textsc{pl}.\textsc{pron}\\
\glt `You could not see any child or any young man at the place where they were butchering turtles.'
\z


\ea\label{ex:text19-33}
A kiu tea abana bebeera koa, paku vonae tea kaku guu mee.\\
\gll a	kiu	te=a	abana	be{$\sim$}beera	koa	paku	vo=na=e	tea	kaku	guu	mee\\
     \textsc{art}2.\textsc{sg}	work	\textsc{prep}=\textsc{art}2.\textsc{sg}	men	\textsc{red}{$\sim$}big	only	do	like=3\textsc{sg}.\textsc{ipfv}=3\textsc{sg}.\textsc{pron}	\textsc{compl}1	butcher	pig	also\\
\glt `It was exclusively the work of the chiefs, it was like butchering pigs.'
\z

\ea\label{ex:text19-34}
Nomana amaa matapaku vai paa rova kahi avuaraara.\\
\gll nomana	a=maa	matapaku	vai	paa	rova	kahi	a=vu=araara\\
     today	\textsc{art}2.\textsc{sg}=\textsc{plm}	practice	\textsc{dem}6	\textsc{tm}3	disappear	from	\textsc{objm}=\textsc{pfv}=1\textsc{pl}.\textsc{in}.\textsc{obj}.\textsc{pron}\\
\glt `Today these practices have disappeared from us.'
\z

\ea\label{ex:text19-35}
Nomana bean tau kaku a goroto -- paku koai; ahiki vira to moroko nae vahaa.\\
\gll nomana	be=an	tau	kaku	a	goroto	paku	koa=i	ahiki	vira	to	moroko	nae	vahaa\\
     today	when=2\textsc{sg}.\textsc{pron}	\textsc{tm}7	butcher	\textsc{art}2.\textsc{sg}	turtle	do	just=3\textsc{sg}.\textsc{pron}	not.exist	indeed	\textsc{nspec}3.\textsc{sg}	talk	3\textsc{sg}.\textsc{poss}.\textsc{pron}	again\\
\glt `Today, when you are about to butcher a turtle -- just do it; indeed, there is not any talk about it.'
\z

\ea\label{ex:text19-36}
AMAA HANANA VAA TEA KAKU GOROTO\\
\gll a=maa	hanana	vaa	tea	kaku	goroto\\
     \textsc{art}2.\textsc{sg}=\textsc{plm}	way	\textsc{lnk}	\textsc{compl}1	butcher	turtle\\
\glt `THE WAYS OF BUTCHERING TURTLES'
\z

\ea\label{ex:text19-37}
(1) Rahi gunaha a goroto teo hum to kahi kaku noman e.\\
\gll rahi	gunaha	a	goroto	te=o	hum	to	kahi	kaku	nom=an	e\\
     pull	go.ashore	\textsc{art}2.\textsc{sg}	turtle	\textsc{prep}=\textsc{art}3.\textsc{sg}	place	\textsc{rel}	\textsc{tm}1	butcher	\textsc{ipfv}=2\textsc{sg}.\textsc{pron}	3\textsc{sg}.\textsc{pron}\\
\glt `(1) Pull the turtle ashore to the place where you are going to butcher it.'
\z

\ea\label{ex:text19-38}
(2) Garuhu vira nio kasuana.\\
\gll garuhu	vira	ni=o	kasuana\\
     wash	thoroughly	\textsc{appl}=\textsc{art}3.\textsc{sg}	sand\\
\glt `(2) Thoroughly wash the sand off.'
\z

\ea\label{ex:text19-39}
(3) Navuhu vamate a goroto, navuhu kii gogoonae.\\
\gll navuhu	va-mate	a	goroto	navuhu	ki=i	gogoo-nae\\
     hit	\textsc{advr}-dead	\textsc{art}2.\textsc{sg}	turtle	hit	\textsc{dat}=3\textsc{sg}.\textsc{pron}	throat-3\textsc{sg}.\textsc{poss}.\textsc{pron}\\
\glt `(3) Hit the turtle dead, hit it on its throat.'
\z

\ea\label{ex:text19-40}
(4) Pee o kapa vaa domonae, pee murina koa o vavaarihi.\\
\gll pee	o	kapa	vaa	domo-nae	pee	murina	koa	o	vavaarihi\\
     cut	\textsc{art}3.\textsc{sg}	shell	\textsc{lnk}	chest-3\textsc{sg}.\textsc{poss}.\textsc{pron}	cut	follow	just	\textsc{art}3.\textsc{sg}	side\\
\glt `Cut the shell on its chest, subsequently cut the sides.'
\z

\ea\label{ex:text19-41}
Ean repaa kapae.\\
\gll e-an	re=paa	kapa=e\\
     \textsc{h}-2\textsc{sg}	\textsc{consec}=\textsc{tm}3	take.off=3\textsc{sg}.\textsc{pron}\\
\glt `Then you take its shell off.'
\z

\ea\label{ex:text19-42}
(5) Rake nao o pusu no hiapa to taneo maana teo orooromo,\\
\gll rake	nao	o	pusu	n=o	hiapa	to	taneo	maa=na	te=o	oro{$\sim$}oromo\\
     Search	thither	\textsc{art}3.\textsc{sg}	upper.end	3\textsc{sg}.\textsc{poss}=\textsc{art}3.\textsc{sg}	intestines	\textsc{rel}	begin	hither=3\textsc{sg}.\textsc{ipfv}	\textsc{prep}=\textsc{art}3.\textsc{sg}	\textsc{red}{$\sim$}swallow\\
\glt `Search for the upper end of the intestines that begin in its throat,'
\z

\ea\label{ex:text19-43}
pee kurukurusui, rahi tavusu kahi bona komana goroto.\\
\gll pee	kuru{$\sim$}kurusu=i	rahi	tavusu	kahi	bona	koma-n=a	goroto\\
     cut	\textsc{red}{$\sim$}cut.off=3\textsc{sg}.\textsc{pron}	pull	come.out	from	\textsc{obj}.\textsc{art}2.\textsc{sg}	inside-3\textsc{sg}.\textsc{poss}=\textsc{art}2.\textsc{sg}	turtle\\
\glt `cut them off, pull them out of the belly of the turtle.'
\z

\ea\label{ex:text19-44}
Rahi tavusu vakavara o hiapa ae amaa meha taba vaa tea goroto,\\
\gll rahi	tavusu	va-kavara	o	hiapa	ae	a=maa	meha	taba	vaa	te=a	goroto\\
     pull	come.out	\textsc{advr}-complete	\textsc{art}3.\textsc{sg}	intestines	and1	\textsc{art}2.\textsc{sg}=\textsc{plm}	other2.\textsc{sg}	thing	\textsc{lnk}	\textsc{prep}=\textsc{art}2.\textsc{sg}	turtle\\
\glt `When you have completely pulled out the instestines and the other things in the turtle,'
\z

\ea\label{ex:text19-45}
(6) kana tavusu o revasin.\\
\gll kana	tavusu	o	revasin\\
     scoop	come.out	\textsc{art}3.\textsc{sg}	blood\\
\glt `scoop the blood out.'
\z

\ea\label{ex:text19-46}
(7) Peepee vihu amaa pakapaka nae (o dorana) bara o kahoo.\\
\gll pee{$\sim$}pee	vihu	a=maa	pakapaka	nae	(o	dorana)	bara	o	kahoo\\
     \textsc{red}{$\sim$}cut	remove	\textsc{art}2.\textsc{plm}	flipper	3\textsc{sg}.\textsc{poss}.\textsc{pron}	\textsc{art}3.\textsc{sg}	four	and2	\textsc{art}3.\textsc{sg}	head\\
\glt `Cut off its flippers (the four) and the head.'
\z

\ea\label{ex:text19-47}
(8) Vaanoto o suraa o beera mene rava bona mata goroto.\\
\gll vaanoto	o	suraa	o	beera	mene	rava	bona	ma-ta	goroto\\
     light	\textsc{art}3.\textsc{sg}	fire	\textsc{art}3.\textsc{sg}	big	for	heat	\textsc{obj}.\textsc{art}2.\textsc{sg}	\textsc{plm}-piece	turtle\\
\glt `Light a big fire for heating the pieces of the turtle.'
\z

\ea\label{ex:text19-48}
Rava vakavara, vaatei a mata goroto teo bana.\\
\gll rava	va-kavara	vaa-tei	a	ma-ta	goroto	te=o	bana\\
     heat	\textsc{advr}-complete	\textsc{caus}-stay	\textsc{art}2.\textsc{sg}	\textsc{plm}-piece	turtle	\textsc{prep}=\textsc{art}3.\textsc{sg}	mat\\
\glt `After having heated them up, put the turtle pieces on a mat.'
\z

\ea\label{ex:text19-49}
(9) Vaagunagunaha a mata goroto, ean repaa peepee vaantee a peepeha mono goroto.\\
\gll vaa-guna{$\sim$}gunaha	a	ma-ta	goroto	e-an	re=paa	pee{$\sim$}pee	va-antee	a	pee{$\sim$}peha	mono	goroto\\
     \textsc{caus}-\textsc{red}{$\sim$}go.down	\textsc{art}2.\textsc{sg}	\textsc{plm}-piece	turtle	\textsc{h}-2\textsc{sg}.\textsc{pron}	\textsc{consec}=\textsc{tm}3	\textsc{red}{$\sim$}cut	\textsc{advr}-enough	\textsc{art}2.\textsc{sg}	\textsc{red}{$\sim$}one	parcel	turtle\\
\glt `Cut the turtle pieces further up, then you cut enough for each turtle parcel.'
\z

\ea\label{ex:text19-50}
(10) Monomono a goroto. Kavara vai voen, vaatei vai tea tao.\\
\gll mono{$\sim$}mono	a	goroto	kavara	vai	vo=en	vaa-tei	vai	te=a	tao\\
     \textsc{red}{$\sim$}wrap.into.parcels	\textsc{art}2.\textsc{sg}	turtle	finished	now	like=\textsc{dem}3	\textsc{caus}-be.there	now	\textsc{prep}=\textsc{art}2.\textsc{sg}	oven\\
\glt `(10) Wrap the turtle up into parcels. When it's finished in this way, put (the parcels) into the earth oven.'
\z

\ea\label{ex:text19-51}
Amaa hum vaa komana goroto na aniani ria ta Teapu.\\
\gll a=maa	hum	vaa	koma-n=a	goroto	na	ani{$\sim$}ani	ri=a	ta	teapu\\
     \textsc{art}2.\textsc{sg}=\textsc{plm}	part	\textsc{lnk}	inside-3\textsc{sg}.\textsc{poss}=\textsc{art}2.\textsc{sg}	turtle	\textsc{tm}2	\textsc{red}{$\sim$}eat	3\textsc{pl}.\textsc{ipfv}=\textsc{art}2.\textsc{sg}	people	Teop\\
\glt `The parts inside the turtle are eaten by the Teop people'
\z

\ea\label{ex:text19-52}
O hiapa na peepee vakakato riori\\
\gll o	hiapa	na	pee{$\sim$}pee	va-kakato	ri=ori\\
     \textsc{art}3.\textsc{sg}	intestines	\textsc{tm}2	\textsc{red}{$\sim$}cut	\textsc{caus}-short	3\textsc{pl}.\textsc{ipfv}=3\textsc{pl}.\textsc{pron}\\
\glt `The intestines are cut into short pieces,'
\z

\ea\label{ex:text19-53}
ore paa siri kopaa bari,\\
\gll o=re	paa	siri	kopaa	bari\\
     3.\textsc{pron}=\textsc{consec}	\textsc{tm}3	tear	cut.lengthwise	4\textsc{pl}.\textsc{pron}\\
\glt `then they are torn into strips,'
\z

\ea\label{ex:text19-54}
ore paa garuhu vamataa bari,\\
\gll o=re	paa	garuhu	va-mataa	bari\\
     3.\textsc{pron}=\textsc{consec}	\textsc{tm}3	wash	\textsc{advr}-good	4\textsc{pl}.\textsc{pron}\\
\glt `then properly washed'
\z

\ea\label{ex:text19-55}
ore paa vaatei bari tea nahu.\\
\gll o=re	paa	vaa-tei	bari	te=a	nahu\\
     3.\textsc{pron}=\textsc{consec}	\textsc{tm}3	\textsc{caus}-be.there	4\textsc{pl}.\textsc{pron}	\textsc{prep}=\textsc{art}2.\textsc{sg}	pot\\
\glt `and put into a pot.'
\z

\ea\label{ex:text19-56}
O gonogonono mee na peepee riori, ore paa vaatei bari tea nahu.\\
\gll o	gonogonono	mee	na	pee{$\sim$}pee	ri=ori	o=re	paa	vaa-tei	bari	te=a	nahu\\
     \textsc{art}3.\textsc{sg}	lungs	also	\textsc{tm}2	\textsc{red}{$\sim$}cut	3\textsc{pl}.\textsc{ipfv}=3\textsc{pl}.\textsc{pron}	3.\textsc{pron}=\textsc{consec}	\textsc{tm}3	\textsc{caus}-be.there	4\textsc{pl}.\textsc{pron}	\textsc{prep}=\textsc{art}2.\textsc{sg}	pot\\
\glt `The lungs are also cut up and put into a pot.'
\z

\ea\label{ex:text19-57}
Amaa hum vai mee, a koburu, o ganusu ae amaa meha hum.\\
\gll a=maa	hum	vai	mee	a	koburu	o	ganusu	ae	a=maa	meha	hum\\
     \textsc{art}2.\textsc{sg}=\textsc{plm}	part	\textsc{dem}6	also	\textsc{art}2.\textsc{sg}	heart	\textsc{art}3.\textsc{sg}	cartilage	and1	\textsc{art}2.\textsc{sg}=\textsc{plm}	other2.\textsc{sg}	part\\
\glt `Also the following parts, the heart, cartilages, and the other parts.'
\z

\ea\label{ex:text19-58}
Be kavara ore paa vaahio bona nahu teo suraa.\\
\gll be	kavara	o=re	paa	vaa-hio	bona	nahu	te=o	suraa\\
     when	finished	3.\textsc{pron}=\textsc{consec}	\textsc{tm}3	\textsc{caus}-sit	\textsc{obj}.\textsc{art}2.\textsc{sg}	pot	\textsc{prep}=\textsc{art}3.\textsc{sg}	fire\\
\glt `When it is finished, the pot is set on the fire.'
\z

\ea\label{ex:text19-59}
Na nahunahu mi riori bono hiihii, ore nin mataa.\\
\gll na	nahu{$\sim$}nahu	mi	ri=ori	bono	hiihii	o=re	nin	mataa\\
     \textsc{tm}2	\textsc{red}{$\sim$}cook	with	3\textsc{pl}.\textsc{ipfv}=3\textsc{pl}.\textsc{pron}	\textsc{obj}.\textsc{art}3.\textsc{sg}	ginger	3.\textsc{pron}=\textsc{consec}	taste	good\\
\glt `It is cooked with ginger, so that it tastes good.'
\z
